
\centering
\def\svgwidth{\columnwidth}
\begin{tikzpicture}[scale=0.8, transform shape, font=\large,
        cell/.style={draw=black, thick, minimum size=1.2cm},
        interior/.style={fill=blue!20},
        global/.style={fill=red!15},
        panel-label/.style={font=\bfseries\large, align=center},
        arrow/.style={->, thick, shorten >=2mm, shorten <=2mm},
        dof/.style={circle, fill=black, inner sep=0.8pt}
    ]

    % Helper to draw patch-wise DoFs (7x7 grid)
    % #1 = top-left cell node name
    % #2 = bottom-right cell node name
    % Requires \usetikzlibrary{calc}
    

    % Helper to draw cubic DoFs for a cell
    % #1 = cell node name
    % Requires \usetikzlibrary{calc}
    

    
    

    % Panel 1: Gather data
    \node[cell] (p1c1) at (0,0) {};
    \node[cell] (p1c2) at (1.2,0) {};
    \node[cell] (p1c3) at (0,-1.2) {};
    \node[cell] (p1c4) at (1.2,-1.2) {};
    \drawInteriorFrame{p1c1}{p1c4}

    \drawCubicDoFs{p1c1}
    \drawCubicDoFs{p1c2}
    \drawCubicDoFs{p1c3}
    \drawCubicDoFs{p1c4}
    \drawPatchInteriorDoFs{p1c1}{p1c4}
    \node[panel-label, above=0.3cm of p1c1.north west, anchor=south west] {Solution\\ Right-hand side\\(all DoFs)};

    % Panel 2: Residual
    \node[cell] (p2c1) at (4.5,0) {};
    \node[cell] (p2c2) at (5.7,0) {};
    \node[cell] (p2c3) at (4.5,-1.2) {};
    \node[cell] (p2c4) at (5.7,-1.2) {};
    \drawInteriorFrame{p2c1}{p2c4}
    \drawPatchInteriorDoFs{p2c1}{p2c4}
    \node[panel-label, above=0.3cm of p2c1.north west, anchor=south west] {Residual \\ (interior DoFs)};

    % Panel 3: Smoothed Residual (moved below Panel 2)
    \node[cell] (p3c1) at (4.5,-3.6) {};
    \node[cell] (p3c2) at (5.7,-3.6) {};
    \node[cell] (p3c3) at (4.5,-4.8) {};
    \node[cell] (p3c4) at (5.7,-4.8) {};
    \drawInteriorFrame{p3c1}{p3c4}
    \drawPatchInteriorDoFs{p3c1}{p3c4}
    \node[panel-label, below=0.3cm of p3c3.south west, anchor=north west] {Residual\\and Correction};

    % Panel 4: Coarse Residual (move between Residual and Correction)
    \node[cell] (p4c1) at (9,-3.6) {};
    \node[cell] (p4c2) at (10.2,-3.6) {};
    \node[cell] (p4c3) at (9,-4.8) {};
    \node[cell] (p4c4) at (10.2,-4.8) {};
    \drawInteriorFrame{p4c1}{p4c4}
    % Show a single central DoF for coarse panels
    \coordinate (p4center) at ($(p4c1.north west)!0.5!(p4c4.south east)$);
    % larger central DoF for coarse panels so it's visible
    % Match interior DoF appearance: blue fill, no outline
    \node[dof, minimum size=10pt, inner sep=0pt, fill=blue, draw=none] at (p4center) {};

    % Panel 5: Coarse Correction (moved below Correction panel)
    \node[cell] (p5c1) at (13.5,-3.6) {};
    \node[cell] (p5c2) at (14.7,-3.6) {};
    \node[cell] (p5c3) at (13.5,-4.8) {};
    \node[cell] (p5c4) at (14.7,-4.8) {};
    \drawInteriorFrame{p5c1}{p5c4}
    % Show a single central DoF for coarse panels
    \coordinate (p5center) at ($(p5c1.north west)!0.5!(p5c4.south east)$);
    % larger central DoF for coarse panels so it's visible
    % Match interior DoF appearance: blue fill, no outline
    \node[dof, minimum size=10pt, inner sep=0pt, fill=blue, draw=none] at (p5center) {};

    % Panel 6: Correction and scatter
    \node[cell] (p6c1) at (13.5,0) {};
    \node[cell] (p6c2) at (14.7,0) {};
    \node[cell] (p6c3) at (13.5,-1.2) {};
    \node[cell] (p6c4) at (14.7,-1.2) {};
    \drawInteriorFrame{p6c1}{p6c4}
    \drawPatchInteriorDoFs{p6c1}{p6c4}
    % bottom-center coordinate for arrows targeting the last panel
    \coordinate (p6bottom) at ($(p6c3.south)!0.5!(p6c4.south)$);
    \node[panel-label, above=0.3cm of p6c1.north west, anchor=south west] {Correction\\(interior only)};

    % Local-solve background: rounded rectangle behind the p2..p6 panels
    % Coordinates with a small padding
    \coordinate (ls_nw) at ($(p2c1.north west)+(-6pt,6pt)$);
    \coordinate (ls_se) at ($(p5c4.south east)+(6pt,-6pt)$);
    % Slightly stronger light-grey fill so the Local Solve region is visible behind the panels
    \draw[rounded corners=6pt, thick, draw=black!50, fill=gray!20, fill opacity=0.40]
    (ls_nw) rectangle (ls_se);
    % Label the local solve box slightly above its top edge; make the label larger and bold
    \node[panel-label, font=\bfseries\LARGE] at ($(ls_nw)!0.5!(ls_se)+(0,0.6cm)$) {Local Solve};

    % Redraw frames, DoFs and labels for panels inside the local-solve box so they appear on top
    % Panel 2: Residual (draw frame without fill so Local Solve background shows through)
    \draw[rounded corners=3pt, thick, draw=blue!60!black] ([shift={(4pt,-4pt)}]p2c1.north west) rectangle
    ([shift={(-4pt,4pt)}]p2c4.south east);
    \drawPatchInteriorDoFs{p2c1}{p2c4}
    \node[panel-label, above=0.3cm of p2c1.north west, anchor=south west] {Residual \\ (interior DoFs)};

    % Panel 3: Smoothed Residual (Residual and Correction)
    \draw[rounded corners=3pt, thick, draw=blue!60!black] ([shift={(4pt,-4pt)}]p3c1.north west) rectangle
    ([shift={(-4pt,4pt)}]p3c4.south east);
    \drawPatchInteriorDoFs{p3c1}{p3c4}
    \node[panel-label, below=0.3cm of p3c3.south west, anchor=north west] {Residual\\and Correction};

    % Panel 4: Coarse Residual
    \draw[rounded corners=3pt, thick, draw=blue!60!black] ([shift={(4pt,-4pt)}]p4c1.north west) rectangle
    ([shift={(-4pt,4pt)}]p4c4.south east);
    \node[panel-label, below=0.3cm of p4c3.south west, anchor=north west] {Coarse\\Residual};

    % Panel 5: Coarse Correction
    \draw[rounded corners=3pt, thick, draw=blue!60!black] ([shift={(4pt,-4pt)}]p5c1.north west) rectangle
    ([shift={(-4pt,4pt)}]p5c4.south east);
    \node[panel-label, below=0.3cm of p5c3.south west, anchor=north west] {Coarse\\Correction};

    % Panel 6: Correction (redraw frame without fill to keep Local Solve background visible)
    \draw[rounded corners=3pt, thick, draw=blue!60!black] ([shift={(4pt,-4pt)}]p6c1.north west) rectangle
    ([shift={(-4pt,4pt)}]p6c4.south east);
    \drawPatchInteriorDoFs{p6c1}{p6c4}
    \node[panel-label, above=0.3cm of p6c1.north west, anchor=south west] {Correction\\(interior only)};

    % Redraw coarse central DoFs on top of the Local Solve box so they are not visually
    % overshadowed by the rounded rectangle background.
    \node[dof, minimum size=10pt, inner sep=0pt, fill=blue, draw=none] at (p4center) {};
    \node[dof, minimum size=10pt, inner sep=0pt, fill=blue, draw=none] at (p5center) {};

    % Arrows between panels
    \draw[arrow] ($(p1c3.west) - (2cm, 0)$) --	node[above, midway, ] {$\overline\Pi_j u, \;\;	\Pi_j b $}
    node[below,
        midway, ] {Gather}	(p1c3.west);
    \draw[arrow] (p1c4.east) -- node[above, midway, ] {$b_j - \Pi_j\overline A_j \overline u_j$} node[below,
        midway, ]
    {Evaluate} (p2c3.west);
    \draw[arrow] (p2c3.south) -- node[midway, right] {Smooth} (p3c1.north);
    % local solve arrow removed (handled implicitly in diagram)
    \draw[arrow] (p6c4.east) -- node[above, midway, ]{$\Pi_j^T d_j$} node[below, midway, ] {Scatter} ($(p6c4.east)
    +
    (2cm, 0)$);

    % V-cycle arrows
    % Make Restrict horizontal at the bottom row (same level as the coarse Smooth)
    \draw[arrow] (p3c4.east) -- node[below, midway] {Restrict} (p4c3.west);
    \draw[arrow] (p4c4.east) -- node[below, midway] {Smooth} (p5c3.west);
    % Make Prolongate vertical from p5 down/up to bottom of last panel (use p6c2.south)
    \draw[arrow] (p5c2.north) -- node[midway, left] {Prolongate and Add} (p6c4.south);

    % Straight dashed arrow from top of last panel to top of second panel (evaluate)
    % Use the same label placement as other arrows (node[above, midway]) so the spacing from the panel is consistent.
    \draw[arrow, dashed] (p6c1.west) --
    node[above, midway] {
        % $b_j - \Pi_j\overline A_j d_j$
    } node[below, midway] {Evaluate} (p2c2.east);

\end{tikzpicture}
